\documentclass[12pt]{article}
\usepackage[utf8]{inputenc}
\usepackage{amsmath}
\usepackage{amsfonts}
\usepackage{amssymb}
\usepackage{enumerate}
\usepackage[margin=1in]{geometry}

\title{GeeksforGeeks Questions}
\author{Generated from JSON}
\date{\today}

\begin{document}

\maketitle

\begin{enumerate}

\item[1.] If A = \{x, y, z\} and B = \{u, v, w, x\}, and the universe is \{s, t, u, v, w, x, y, z\}. Then (A U B') ∩ (A ∩ B) is equal to
\begin{enumerate}[label=\Alph*.]
  \item \{u, v, w, x\}
\end{enumerate}

\item[2.] Minimum number of NAND gates required to implement the following binary equation (A'+B')(C+D),
\begin{enumerate}[label=\Alph*.]
  \item 4
  \item 5
  \item 3
  \item 6
\end{enumerate}

\item[3.] Of the following, which best approximates the ratio of the number of non-terminal nodes in the total number of nodes in a complete K-ary tree of depth N ?
\begin{enumerate}[label=\Alph*.]
  \item 1/N
  \item (N-1)/N
  \item 1/K
  \item (K-1)/K
\end{enumerate}

\item[4.] In a class definition with 10 methods, to make the class maximally cohesive, number of direct and indirect connections required among the methods are
\begin{enumerate}[label=\Alph*.]
  \item 90, 0
  \item 45, 0
  \item 10, 10
  \item 45, 45
\end{enumerate}

\item[5.] Consider the following pseudo-code:
\begin{enumerate}[label=\Alph*.]
  \item 3
  \item 2
  \item 4
  \item 1
\end{enumerate}

\item[6.] Raymonds tree based algorithm ensures
\begin{enumerate}[label=\Alph*.]
  \item no starvation, but deadlock may occur in rare cases
  \item no deadlock, but starvation may occur
  \item neither deadlock nor starvation can occur
  \item deadlock may occur in cases where the process is already starved
\end{enumerate}

\item[7.] The Master theorem
\begin{enumerate}[label=\Alph*.]
  \item assumes the subproblems are unequal sizes
  \item can be used if the subproblems are of equal size
  \item cannot be used for divide and conquer algorithms
  \item cannot be used for asymptotic complexity analysis
\end{enumerate}

\item[8.] The minimum height of an AVL tree with n nodes is
\begin{enumerate}[label=\Alph*.]
  \item Ceil (log
  \item (n + 1))
  \item 1.44 log
  \item n
\end{enumerate}

\item[9.] What is the in-order successor of 15 in the given binary search tree ?
\begin{enumerate}[label=\Alph*.]
  \item 18
  \item 6
  \item 17
  \item 2
\end{enumerate}

\item[10.] Consider the following circuit
\begin{enumerate}[label=\Alph*.]
  \item (AB)'E + EF + (CD)'F
  \item (E' + ABF')(C + D + F')
  \item ((AB)' + E)(E' + F')(C + D + F')
  \item (A + B)E' + (EF)' + CDF'
\end{enumerate}

\item[11.] If ABCD is a 4-bit binary number, then what is code generated by the following circuit ?
\begin{enumerate}[label=\Alph*.]
  \item BCD code
  \item Gray code
  \item 8421 code
  \item Excess-3 code
\end{enumerate}

\item[12.] A grammar is defined as
\begin{enumerate}[label=\Alph*.]
  \item \{A, B, C, D, E\}
  \item \{B, C D, E\}
  \item \{A, B, C, D, E, x, y, z\}
  \item \{x, y, z\}
\end{enumerate}

\item[13.] The operating system and the other processes are protected from being modified by an already running process because
\begin{enumerate}[label=\Alph*.]
  \item they run at different time instants and not in parallel
  \item they are in different logical addresses
  \item they use a protection algorithm in the scheduler
  \item every address generated by the CPU is being checked against the relocation and limit parameters
\end{enumerate}

\item[14.] What is compaction refers to
\begin{enumerate}[label=\Alph*.]
  \item a technique for overcoming internal fragmentation
  \item a paging technique
  \item a technique for overcoming external fragmentation
  \item a technique for compressing the data
\end{enumerate}

\item[15.] In linear hashing, if blocking factor bfr, loading factor i and file buckets N are known, the number of records will be
\begin{enumerate}[label=\Alph*.]
  \item cr = i + bfr + N
  \item r = i – bfr – N
  \item r = i – bfr – N
  \item r = i * bfr * N
\end{enumerate}

\item[16.] The post-order traversal of a binary tree is ACEDBHIGF. The pre-order traversal is
\begin{enumerate}[label=\Alph*.]
  \item ABCDEFGHI
  \item FBADCEGIH
  \item FABCDEGHI
  \item ABDCEFGIH
\end{enumerate}

\item[17.] Three CPU-bound tasks, with execution times of 15, 12 and 5 time units respectively arrive at times 0, t and 8, respectively. If the operating system implements a shortest remaining time first scheduling algorithm, what should be the value of t to have 4 context switches ? Ignore the context switches at time 0 and at the end.
\begin{enumerate}[label=\Alph*.]
  \item 0 \textless{} t \textless{} 3
  \item t = 0
  \item t \textless{} = 3
  \item 3 \textless{} t \textless{} 8
\end{enumerate}

\item[18.] Consider the following page reference string.
\begin{enumerate}[label=\Alph*.]
  \item 7
  \item 4
  \item 6
  \item 5
\end{enumerate}

\item[19.] An aid to determine the deadlock occurrence is
\begin{enumerate}[label=\Alph*.]
  \item resource allocation graph
  \item starvation graph
  \item inversion graph
  \item none of the above
\end{enumerate}

\item[20.] Dispatch latency is defined as
\begin{enumerate}[label=\Alph*.]
  \item the speed of dispatching a process from running to the ready state
  \item the time of dispatching a process from running to ready state and keeping the CPU idle
  \item the time to stop one process and start running another one
  \item none of these
\end{enumerate}

\item[21.] Context free languages are closed under
\begin{enumerate}[label=\Alph*.]
  \item union, intersection
  \item union, kleene closure
  \item intersection, complement
  \item complement, kleene closure
\end{enumerate}

\item[22.] Consider a 2-dimensional array x with 10 rows and 4 columns, with each element storing a value equivalent to the product of row number and column number. The array row major format. If the first element x[0][0] occupies the memory location with address 1000 and each element occupies only one memory location, which all locations (in decimal) will be holding a value of 10 ?
\begin{enumerate}[label=\Alph*.]
  \item 1018, 1019
  \item 1022, 1041
  \item 1013, 1014
  \item 1000, 1399
\end{enumerate}

\item[23.] Which one indicates a Technics of building cross compilers ?
\begin{enumerate}[label=\Alph*.]
  \item Beta cross
  \item Canadian cross
  \item Mexican cross
  \item X-cross
\end{enumerate}

\item[24.] A computer which issues instructions in order, has only 2 registers and 3 opcodes ADD, SUB and MOV. Consider 2 different implementations of the following basic block:
\begin{enumerate}[label=\Alph*.]
  \item Case 2, 2
  \item Case 2, 3
  \item Case 1, 2
  \item Case 1, 3
\end{enumerate}

\item[25.] A stack organised computer is characterised by instructions with
\begin{enumerate}[label=\Alph*.]
  \item indirect addressing
  \item direct addressing
  \item zero addressing
  \item index addressing
\end{enumerate}

\item[26.] Which of the following is a type of a out-of-order execution, with the reordering done by a compiler
\begin{enumerate}[label=\Alph*.]
  \item loop unrolling
  \item dead code elimination
  \item strength reduction
  \item software pipelining
\end{enumerate}

\item[27.] Minimum number of states required in DFA accepting binary strings not ending in “101” is ,
\begin{enumerate}[label=\Alph*.]
  \item 3
  \item 4
  \item 5
  \item 6
\end{enumerate}

\item[28.] Which of the following classes of languages can validate an IPv4 address in dotted decimal format? It is to be ensured that the decimal values are between 0 and 255
\begin{enumerate}[label=\Alph*.]
  \item RE and higher
  \item CFG and higher
  \item CSG and higher
  \item Recursively enumerable language
\end{enumerate}

\item[29.] The language which is generated by the grammar S → aSa\textbar{}bSb\textbar{}a\textbar{}b over the alphabet \{a, b\} is the set of
\begin{enumerate}[label=\Alph*.]
  \item Strings that begin and end with the same symbol
  \item All odd and even length palindromes
  \item All odd length palindromes
  \item All even length palindromes
\end{enumerate}

\item[30.] Which of the following is true ?
\begin{enumerate}[label=\Alph*.]
  \item Every subset of a regular set is regular
  \item Every finite subset of non-regular set is regular
  \item The union of two non-regular set is not regular
  \item Infinite union of finite set is regular
\end{enumerate}

\item[31.] Which of the following algorithms defines time quantum ?
\begin{enumerate}[label=\Alph*.]
  \item shortest job scheduling algorithm
  \item round robin scheduling algorithm
  \item priority scheduling algorithm
  \item multilevel queue scheduling algorithm
\end{enumerate}

\item[32.] A non-pipelined CPU has 12 general purpose registers (R0, R1, R2, ..... R12). Following operation are supported:
\begin{enumerate}[label=\Alph*.]
  \item 5
  \item 6
  \item 7
  \item 8
\end{enumerate}

\item[33.] An array of 2 two byte integers is stored in big endian machine in byte address as shown below. What will be its storage pattern in little endian machine ?
\begin{enumerate}[label=\Alph*.]
  \item 0 × 104 12
  \item 0 × 104 12
  \item 0 × 103 56
  \item 0 × 102 34
\end{enumerate}

\item[34.] The persist timer is used in TCP to
\begin{enumerate}[label=\Alph*.]
  \item To detect crashes from the other end of the connection
  \item To enable retransmission
  \item To avoid deadlock condition
  \item To timeout FIN\_Wait1 condition
\end{enumerate}

\item[35.] To send same bit sequence, NRZ encoding require
\begin{enumerate}[label=\Alph*.]
  \item Same clock frequency as Manchester encoding
  \item Half the clock frequency as Manchester encoding
  \item Twice the clock frequency as Manchester encoding
  \item A clock frequency which depend on number of zeros and ones in the bit sequence
\end{enumerate}

\item[36.] The SQL query
\begin{enumerate}[label=\Alph*.]
  \item All rows in Table B, which meets equality condition above and, none from Table A, which meets the condition
  \item All rows in Table A, which meets equality condition above and none from Table B, which meets the condition
  \item All rows in Table B, which meets equality condition
  \item All rows in Table A, which meets equality condition
\end{enumerate}

\item[37.] If every non-key attribute functionally dependent on the primary key, then the relation will be in
\begin{enumerate}[label=\Alph*.]
  \item First normal form
  \item Second normal form
  \item Third normal form
  \item Fourth normal form
\end{enumerate}

\item[38.] One instruction tries to write an operand before it is written by previous instruction. This may lead to a dependency called
\begin{enumerate}[label=\Alph*.]
  \item True dependency
  \item Anti-dependency
  \item Output dependency
  \item Control hazard
\end{enumerate}

\item[39.] In a two-pass assembler, resolution of subroutine calls and inclusion of labels in the symbol table is done during
\begin{enumerate}[label=\Alph*.]
  \item second pass
  \item first pass and second pass respectively
  \item second pass and first pass respectively
  \item first pass
\end{enumerate}

\item[40.] The number of tokens in the following C code segment is
\begin{enumerate}[label=\Alph*.]
  \item 27
  \item 29
  \item 26
  \item 24
\end{enumerate}

\item[41.] Given that
\begin{enumerate}[label=\Alph*.]
  \item Every fish is eaten by some bear
  \item Bears eat only fish
  \item Every bear eats fish
  \item Only bears eat fish
\end{enumerate}

\item[42.] Given the grammar:
\begin{enumerate}[label=\Alph*.]
  \item Grammar is not ambiguous
  \item Priority of + over * is ensured
  \item Right to left evaluation of * and + happens
  \item None of these
\end{enumerate}

\item[43.] Huffman tree is constructed for the following data: \{A, B, C, D, E\} with frequency \{0.17, 0.11, 0.24, 0.33 and 0.15\} respectively. 100 00 01101 is decoded as
\begin{enumerate}[label=\Alph*.]
  \item BACE
  \item CADE
  \item BAD
  \item CADD
\end{enumerate}

\item[44.] If the array A contains the items 10, 4, 7, 23, 67, 12 and 5 in that order, what will be the resultant array A after third pass of insertion sort ?
\begin{enumerate}[label=\Alph*.]
  \item 67, 12, 10, 5, 4, 7, 23
  \item 4, 7, 10, 23, 67, 12, 5
  \item 4, 5, 7, 67, 10, 12, 23
  \item 10, 7, 4, 67, 23, 12, 5
\end{enumerate}

\item[45.] G is an undirected graph with vertex set \{v1, v2, v3, v4, v5, v6, v7\} and edge set \{v1v2, v1v3, v1v4, v2v4, v2v5, v3v4, v4v5, v4v6, v5v6, v6v7\}. A breadth first search of the graph is performed with v1 as the root node. Which of the following is a tree edge ?
\begin{enumerate}[label=\Alph*.]
  \item v2v4
  \item v1v4
  \item v4v5
  \item v3v4
\end{enumerate}

\item[46.] Convert the pre-fix expression to in-fix
\begin{enumerate}[label=\Alph*.]
  \item None
  \item (A – B)*C + (D*E) – (F + G)
  \item (A + B)*C – (D – E)*(F – G)
  \item (A + B – C)*(D – E))*(F + G)
\end{enumerate}

\item[47.] Which of the following affects the processing power assuming they do not influence each other.
\begin{enumerate}[label=\Alph*.]
  \item 3 only
  \item 1 and 3 only
  \item 2 and 3 only
  \item 1, 2 and 3
\end{enumerate}

\item[48.] Statements associated with registers of a CPU are given. Identify the false statement.
\begin{enumerate}[label=\Alph*.]
  \item The program counter holds the memory address of the instruction in execution.
  \item Only opcode is transferred to the control unit.
  \item An instruction in the instruction register consists of the opcode and the operand.
  \item The value of the program counter is incremented by 1 once its value has been read to the memory address register.
\end{enumerate}

\item[49.] The immediate addressing mode can be used for
\begin{enumerate}[label=\Alph*.]
  \item Only 1
  \item Only 2
  \item Both 1 and 2
  \item Immediate mode refers to data in cache
\end{enumerate}

\item[50.] Following declaration of an array of struct, assumes size of byte, short, int and long are 1, 2, 3 and 4 respectively. Alignment rule stipulates that n-byte field must be located at an address divisible by n. The fields in a struct are not rearranged, padding is used to ensure alignment. All elements of array should be of same size.
\begin{enumerate}[label=\Alph*.]
  \item 150
  \item 160
  \item 200
  \item 240
\end{enumerate}

\item[51.] What is the complexity of the following code ?
\begin{enumerate}[label=\Alph*.]
  \item Ο(n
  \item )
  \item Ο(n log n)
  \item Ο(n)
\end{enumerate}

\item[52.] A stack is implemented with an array of ‘A[0...N – 1]’ and a variable ‘pos’. The push and pop operations are defined by the following code.
\begin{enumerate}[label=\Alph*.]
  \item pos ← –1
  \item pos ← 0
  \item pos ← 1
  \item pos ← N – 1
\end{enumerate}

\item[53.] Following Multiplexer circuit is equivalent to
\begin{enumerate}[label=\Alph*.]
  \item Sum equation of full adder
  \item Carry equation of full adder
  \item Borrow equation for full subtractor
  \item Difference equation of a full subtractor
\end{enumerate}

\item[54.] In a 8-bit ripple carry adder using identical full adders, each full adder takes 34 ns for computing sum. If the time taken for 8-bit addition is 90 ns, find time taken by each full adder to find carry.
\begin{enumerate}[label=\Alph*.]
  \item 6 ns
  \item 7 ns
  \item 10 ns
  \item 8 ns
\end{enumerate}

\item[55.] Consider a 32-bit processor which supports 70 instructions. Each instruction is 32 bit long and has 4 fields namely opcode, two register identifiers and an immediate operand of unsigned integer type. Maximum value of the immediate operand that can be supported by the processor is 8191. How many registers the processor has ?
\begin{enumerate}[label=\Alph*.]
  \item 32
  \item 64
  \item 128
  \item 16
\end{enumerate}

\item[56.] Consider a 5-segment pipeline with a clock cycle time 20 ns in each sub operation. Find out the approximate speed-up ratio between pipelined and non-pipelined system to execute 100 instructions. (If an average, every five cycles, a bubble due to data hazard has to be introduced in the pipeline.).
\begin{enumerate}[label=\Alph*.]
  \item 5
  \item 4.03
  \item 4.81
  \item 4.17
\end{enumerate}

\item[57.] What is the defect rate for Six sigma ?
\begin{enumerate}[label=\Alph*.]
  \item 1.0 defect per million lines of code
  \item 1.4 defects per million lines of code
  \item 3.0 defects per million lines of code
  \item 3.4 defects per million lines of code
\end{enumerate}

\item[58.] What is the availability of the software with following reliability figures.
\begin{enumerate}[label=\Alph*.]
  \item 90\%
  \item 96\%
  \item 24\%
  \item 50\%
\end{enumerate}

\item[59.] The following circuit compares two 2-bit binary numbers, X and Y represented by X1X0 and Y1Y0 respectively. (X0 and Y0 represent Least Significant Bits)
\begin{enumerate}[label=\Alph*.]
  \item X \textgreater{} Y
  \item X \textless{} Y
  \item X = Y
  \item X! = Y
\end{enumerate}

\item[60.] Of the following sort algorithms, which has execution time that is least dependent on initial ordering of the input ?
\begin{enumerate}[label=\Alph*.]
  \item Insertion sort
  \item Quick sort
  \item Merge sort
  \item Selection sort
\end{enumerate}

\item[61.] How many total bits are required for a direct-mapped cache with 128 KB of data and 1 word block size, assuming a 32-bit address and 1 word size of 4 bytes ?
\begin{enumerate}[label=\Alph*.]
  \item 2 Mbits
  \item 1.7 Mbits
  \item 2.5 Mbits
  \item 1.5 Mbits
\end{enumerate}

\item[62.] The hardware implementation which provides mutual exclusion is
\begin{enumerate}[label=\Alph*.]
  \item Semaphores
  \item Test and set instruction
  \item Both options
  \item None of the options
\end{enumerate}

\item[63.] For the distributions given below :
\begin{enumerate}[label=\Alph*.]
  \item Standard deviation of A is significantly lower than standard deviation of B
  \item Standard deviation of A is slightly lower than standard deviation of B
  \item Standard deviation of A is same as standard deviation of B
  \item Standard deviation of A is significantly higher than standard deviation of B
\end{enumerate}

\item[64.] If x + 2y = 30, then (2y/5 + x/3) + (x/5 + 2y/3) will be equal to
\begin{enumerate}[label=\Alph*.]
  \item 8
  \item 16
  \item 18
  \item 20
\end{enumerate}

\item[65.] Checksum field in TCP header is
\begin{enumerate}[label=\Alph*.]
  \item ones complement of sum of header and data in bytes
  \item ones complement of sum of header, data and pseudo header in 16 bit words
  \item dropped from IPv6 header format
  \item better than md5 or sh1 methods
\end{enumerate}

\item[66.] What is the output of the code given below?
\begin{enumerate}[label=\Alph*.]
  \item 100
  \item 110
  \item 40
  \item 44
\end{enumerate}

\item[67.] A given grammar is called ambiguous if ,
\begin{enumerate}[label=\Alph*.]
  \item two or more productions have the same non-terminal on the left hand side
  \item a derivation tree has more than one associated sentence
  \item there is a sentence with more than one derivation tree corresponding to it
  \item brackets are not present in the grammar
\end{enumerate}

\item[68.] Consider the following recursive C function that takes two arguments
\begin{enumerate}[label=\Alph*.]
  \item 9
  \item 8
  \item 5
  \item 2
\end{enumerate}

\item[69.] Remote Procedure Calls are used for
\begin{enumerate}[label=\Alph*.]
  \item communication between two processes remotely different from each other on the same system.
  \item communication between two processes on the same system.
  \item communication between two processes on separate system.
  \item None of the above.
\end{enumerate}

\item[70.] Properties of ‘DELETE’ and ‘TRUNCATE’ commands indicate that
\begin{enumerate}[label=\Alph*.]
  \item After the execution of ‘TRUNCATE’ operation, COMMIT and ROLLBACK statements cannot be performed to retrieve the lost data, while ‘DELETE’ allow it.
  \item After the execution of ‘DELETE’ and ‘TRUNCATE’ operation retrieval is easily possible for the lost data.
  \item After the execution of ‘DELETE’ operation, COMMIT and ROLLBACK statements can be performed to retrieve the lost data, while TRUNCATE do not allow it.
  \item After the execution of ‘DELETE’ and ‘TRUNCATE’ operation no retrieval is possible for the lost data.
\end{enumerate}

\item[71.] Regression testing is primarily related to
\begin{enumerate}[label=\Alph*.]
  \item Functional testing
  \item Development testing
  \item Data flow testing
  \item Maintenance testing
\end{enumerate}

\item[72.] A magnetic disk has 100 cylinders, each with 10 tracks of 10 sectors. If each sector contains 128 Bytes, what is the maximum capacity of the disk in kilobytes ?
\begin{enumerate}[label=\Alph*.]
  \item 1,280,000
  \item 1280
  \item 1250
  \item 128,000
\end{enumerate}

\item[73.] Avalanche effect in cryptography refers
\begin{enumerate}[label=\Alph*.]
  \item Large changes in cipher text when the keyword is changed minimally
  \item Large changes in cipher text when the plain text is changed
  \item Large impact of keyword change to length of the cipher text
  \item None of the above
\end{enumerate}

\item[74.] In a columnar transposition cipher, the plain text is “the tomato is a plant in the night shade family”, keyword is “TOMATO”.   The cipher text is
\begin{enumerate}[label=\Alph*.]
  \item “TINESAX / EOAHTFX / HTLTHEY / MAIIAIX / TAPNGDL / OSTNHMX”
  \item “TINESAX / EOAHTFX / MAIIAIX / HTLTHEY / TAPNGDL / OSTNHMX”
  \item “TINESAX / EOAHTFX / HTLTHEY / MAIIAIX / OSTNHMX / TAPNGDL”
  \item “EOAHTFX / TINESAX / HTLTHEY / MAIIAIX / TAPNGDL / OSTNHMX”
\end{enumerate}

\item[75.] Which of the following is an efficient method of cache updating ?
\begin{enumerate}[label=\Alph*.]
  \item Snoopy writes
  \item Write through
  \item Write within
  \item Buffered write
\end{enumerate}

\item[76.] What is the output in a 32 bit machine with 32 bit compiler ?
\begin{enumerate}[label=\Alph*.]
  \item 60 70
  \item 50 50
  \item 50 60
  \item 60 50
\end{enumerate}

\item[77.] What is the output of the following ‘c’ code assuming it runs on a byte addressed little endian machine?
\begin{enumerate}[label=\Alph*.]
  \item 622
  \item 311
  \item 22
  \item 110
\end{enumerate}

\item[78.] A new flip flop with inputs X and Y, has the following property.
\begin{enumerate}[label=\Alph*.]
  \item X'Q' + Y'Q
  \item X'Q + Y'Q'
  \item XQ' + YQ
  \item XQ' + Y'Q
\end{enumerate}

\item[79.] Consider product of three matrices M1, M2 and M3 having w rows and x columns, x rows and y columns, and y rows and z columns. Under what condition will it take less time to compute the product as (M1M2)M3 than to compute M1(M2M3) ?
\begin{enumerate}[label=\Alph*.]
  \item Always take the same time
  \item (1/x + 1/z) \textless{} (1/w + 1/y)
  \item x \textgreater{} y
  \item (w + x) \textgreater{} (y + z)
\end{enumerate}

\item[80.] In the following procedure
\begin{enumerate}[label=\Alph*.]
  \item 5
  \item 8
  \item 13
  \item 0
\end{enumerate}

\end{enumerate}

\end{document}